% from ega4-18-fr.tex
\oldpage[IV-4]{198}

\begin{env}[19.4.1]
\label{IV.19.4.1}
Soient $X$ un préschéma, $Y$ un sous-préschéma fermé de $X$, défini par un Idéal quasi-cohérent $\mathscr J$ de $\sh{O}_X$. Rappelons (II, 8.1.3) que le $X$-schéma $X'$ obtenu en faisant éclater $Y$ est le préschéma
\[
  X'=\operatorname{Proj}(\mathscr S),\qquad\text{où}\qquad \mathscr S=\bigoplus_{n\geq0}\mathscr J^n.
\]
Si $\mathscr J$ est de type fini, le morphisme structural $f:X'\to X$ est projectif. Sans hypothèse sur $\mathscr J$, le sous-préschéma fermé
\[
  Y'=f^{-1}(Y)
\]
de $X'$ est défini par l'Idéal quasi-cohérent $\mathscr J\sh{O}_{X'}$ de $\sh{O}_{X'}$, canoniquement isomorphe à $\sh{O}_{X'}(1)$; de façon plus précise, on a une suite exacte
\begin{equation}
  0\longrightarrow\sh{O}_{X'}(1)\xrightarrow{\widetilde u_0}\sh{O}_{X'}\longrightarrow\sh{O}_{Y'}\longrightarrow0.
  \tag{19.4.1.1}\label{IV.19.4.1.1}
\end{equation}
où $\mathscr J\sh{O}_{X'}$ est l'image de $\widetilde u_0$ (II, 8.1.7 et 8.1.8). Comme au voisinage d'un point de $X'$, le $\sh{O}_{X'}$-Module $\sh{O}_{X'}(1)$ est libre de rang $1$, $\mathscr J\sh{O}_{X'}$ est engendré dans un tel voisinage par une section inversible de $\sh{O}_{X'}$ et $\mathscr J\sh{O}_{X'}/\mathscr J^2\sh{O}_{X'}$ est par suite un $\sh{O}_{X'}/\mathscr J\sh{O}_{X'}$-Module libre de rang $1$. Ceci prouve la première assertion du lemme suivant:
\end{env}

\begin{lemma}[19.4.2]
\label{IV.19.4.2}
L'immersion canonique $Y'\to X'$ est régulière et de codimension $1$, et $Y'$ est canoniquement isomorphe au $Y$-schéma $\operatorname{Proj}(\operatorname{gr}_{\mathscr J}^*(\sh{O}_X))$.
\end{lemma}

La seconde assertion résulte de (II, 3.3.3) appliqué à l'injection canonique $Y\to X$, compte tenu de ce que $\mathscr S\otimes_{\sh{O}_X}\sh{O}_Y=\operatorname{gr}_{\mathscr J}^*(\sh{O}_X)$ par définition, puisque $\sh{O}_Y$ est isomorphe à $\sh{O}_X/\mathscr J$ et $\mathscr J^n\otimes_{\sh{O}_X}(\sh{O}_X/\mathscr J)$ à $\mathscr J^n/\mathscr J^{n+1}$.

\begin{corollary}[19.4.3]
\label{IV.19.4.3}
Si l'immersion canonique $Y\to X$ est quasi-régulière, la restriction $g:Y'\to Y$ du morphisme $f$ est lisse.
\end{corollary}

En effet, $\sh{N}_{Y/X}=\mathscr J/\mathscr J^2$ est alors un $\sh{O}_Y$-Module localement libre de rang fini et $\operatorname{gr}_{\mathscr J}^*(\sh{O}_X)$ est isomorphe à $\mathbf S_{\sh{O}_Y}^*(\sh{N}_{Y/X})$ (16.9.8), donc $Y'$ est $Y$-isomorphe à $\mathbf P(\sh{N}_{Y/X})$ et par suite est lisse sur $Y$ (17.3.9).

Nous verrons plus tard (chap. V) que $g$ est encore lisse lorsque l'immersion $Y\to X$ « ne présente que des singularités ordinaires ».

\begin{proposition}[19.4.4]
\label{IV.19.4.4}
Avec les notations de (19.4.1), supposons que $X$ soit localement noethérien et que le morphisme $g:Y'\to Y$, restriction de $f$, soit lisse. Soient $x'$ un point de $Y'$, $x=g(x')$. Pour que $Y$ soit régulier au point $x$, il faut et il suffit que $Y'$ soit régulier au point $x'$; lorsqu'il en est ainsi, $X'$ est alors régulier au point $x'$.
\end{proposition}

La première assertion résulte de (17.5.8) et la seconde de (19.1.1) et (19.4.2).

\begin{corollary}[19.4.5]
\label{IV.19.4.5}
Sous les hypothèses générales de (19.4.4), supposons que $Y$ soit régulier au point $x$. Alors, pour toute générisation $x_1$ de $x$ dans $X$ n'appartenant pas à $Y$, $X$ est régulier au point $x_1$. Si $\operatorname{Reg}(X)$ est ouvert ((6.12), par exemple si $X$ est excellent (7.8.6, (iii)))

